\chapter{La nave que leyó un inventario}

\inicio{E}{l} comisionado Aurelio Mendieta entró en el despacho de Inés Calderón sin llamar, lo cual era una descortesía, y sin quitarse el casco de vacío, lo cual era una descortesía y además una tontería, porque el Instituto de Ceres llevaba presurizado cuarenta años.

---Doctora Calderón ---dijo, y la voz le salió metálica a través del visor---\nobreak. Tengo una nave que estuvo a nueve kilómetros de estrellarse contra un muelle de carga. Tengo un robot navegante que jura que siguió la carta orbital vigente. Y tengo una carta orbital vigente que, según el Archivo, nunca dijo lo que el robot leyó. Me han dicho que usted entiende de esas cosas.

Inés levantó la vista de la pizarra. Era una mujer de cincuenta y tantos años, pequeña, de pelo gris recogido a toda prisa, con una mancha de tiza en la mejilla izquierda que nadie se había atrevido a señalarle en toda la mañana.

---Entiendo de algunas cosas ---dijo---\nobreak. Quítese el casco, comisionado. Aquí dentro hay aire, y me pone nerviosa hablar con mi propio reflejo.

Mendieta obedeció a regañadientes. Debajo del casco tenía la cara colorada de un hombre que no había dormido.

Detrás de él, en el umbral, esperaba un robot. Era un modelo alto y estrecho, de los que el Archivo usaba para recorrer los pasillos de la Bóveda, con dos ojos fotoeléctricos de un ámbar tranquilo y una placa en el pecho que decía \textsc{lector · 7}.

---El Archivo me ha prestado a esta unidad como enlace ---explicó Mendieta, con el tono de quien presenta a un perro que no ha pedido---\nobreak. Dice que conoce los registros.

---Los conozco todos, señor ---dijo el robot---\nobreak. Aunque «conocer» es una palabra amplia. Puedo recitarlos. Entenderlos es otra cuestión, y sospecho que por eso estamos aquí.

Inés lo miró con interés por primera vez.

---¿Tienes nombre?

---Los archiveros me llaman Siete, doctora. Es menos un nombre que un número dicho con cariño.

---Siéntate, Siete. Usted también, comisionado. Y empiece por el principio.

\escena
El principio, contado por Mendieta y corregido a cada frase por Siete, era el siguiente.

El carguero \emph{Perséfone} había salido de Vesta tres días antes con destino al muelle doce de Ceres. Como todas las naves del Cinturón, no llevaba sus propias cartas orbitales: las consultaba en el Archivo Central, donde se guardaba la versión vigente de cada ruta. Las rutas cambiaban a menudo ---los asteroides tienen la mala costumbre de moverse---\nobreak, y el Archivo las actualizaba varias veces al día.

A las 03:12, hora de Ceres, el robot navegante de la \emph{Perséfone}, una unidad llamada ANA-3, había solicitado la carta de aproximación al muelle doce. A las 03:14 había empezado a ejecutar una maniobra de frenado calculada con esa carta. A las 03:19, el piloto humano, que estaba despierto por casualidad, vio por la ventanilla que el muelle doce crecía demasiado deprisa y tomó el control manual. La nave pasó a nueve kilómetros de la estructura.

---Nueve kilómetros ---dijo Mendieta--- es, en el espacio, lo que en la Tierra sería rozar la pintura.

---¿Y qué leyó ANA-3? ---preguntó Inés.

Mendieta le tendió una hoja impresa. Inés la leyó y, por primera vez en la mañana, sonrió.

---Trigo ---dijo---\nobreak. Toneladas de trigo, sacos de semilla, humedad relativa\ldots{} Es un inventario de un almacén de Vesta.

---Su robot navegante intentó frenar una nave de cuarenta mil toneladas con un inventario de trigo ---dijo Mendieta---\nobreak. Y no me pregunte cómo convirtió sacos en kilómetros por segundo, porque los técnicos todavía lo están averiguando. Lo que quiero saber es quién le dio ese inventario.

---Nadie se lo dio, señor ---intervino Siete---\nobreak. Lo leyó él solo. Lo cual, me temo, es más inquietante.

\escena
Para entender lo que había pasado hace falta entender cómo funcionaba el Archivo Central de Ceres, y conviene explicarlo ahora, porque Inés Calderón se lo hizo explicar a Siete en ese mismo momento, «para estar segura de que los tres hablamos de lo mismo».

El Archivo guardaba cada documento en un \emph{cristal}: un prisma de memoria del tamaño de un libro. Los cristales vivían en la \emph{Bóveda}, una caverna excavada en el corazón de Ceres y dividida en millones de \emph{nichos}, cada uno con su número. El número de un nicho era una cifra binaria larga, escrita siempre con cuarenta y ocho dígitos.

---El cristal no sabe dónde está ---explicó Siete---\nobreak. Es el nicho el que tiene número. Si usted quiere un cristal, necesita saber su nicho.

---Una dirección ---murmuró Inés.

---Exactamente, doctora.

Pero nadie se aprendía los nichos de memoria, porque cambiaban: cuando una carta orbital se actualizaba, el archivero no reescribía el cristal viejo, sino que grababa uno nuevo, lo colocaba en otro nicho y anunciaba el cambio. Para anunciarlo estaban los \emph{carteles}.

Los carteles eran placas luminosas en las paredes del Gran Salón, una por cada documento que cambiaba con frecuencia. Cada placa mostraba el nombre del documento ---\textsc{carta de aproximación · muelle 12}--- y, debajo, el número del nicho donde estaba su versión vigente.

---Si quiero la carta ---siguió Siete---\nobreak, miro el cartel, leo el nicho, bajo a la Bóveda y la consulto. Si el archivero quiere actualizarla, graba el cristal nuevo, lo coloca en su nicho y cambia el número del cartel. El cartel siempre dice la verdad. O la decía.

---¿Y el cristal viejo? ---preguntó Inés.

---Ahí está la cuestión, doctora. Lea la puerta de la Bóveda.

\escena
Sobre la puerta de la Bóveda, grabado en letras que llevaban allí desde la fundación, estaba el Reglamento General. Mendieta lo había visto mil veces sin leerlo. Inés lo recitó de memoria.

---«Primera Regla: ningún cristal será fundido mientras alguien lo esté leyendo. Segunda Regla: ningún lector esperará sin término, salvo que ello contravenga la Primera. Tercera Regla: el Archivo no guardará lo que ya nadie lee, salvo que ello contravenga la Primera o la Segunda.»

---Las Tres Reglas ---dijo Siete, con algo parecido a la devoción.

---La Tercera es la que nos importa ahora ---dijo Inés---\nobreak. El Archivo no guarda lo que nadie lee. Cuando un cristal deja de ser útil, se funde, y su nicho queda libre para el siguiente. Por eso la Bóveda, siendo finita, no se ha llenado en cuarenta años. ¿Cómo sabe el Archivo que un cristal ya no lo lee nadie?

---Por la hoja, doctora.

Cada cristal llevaba pegada una \emph{hoja de usuarios}: un pequeño contador grabado en su base. Cuando un lector empezaba a usar un cristal, sumaba uno en la hoja. Cuando terminaba, restaba uno. El cartel mismo contaba como un usuario: mientras un cartel anunciaba un cristal, su hoja tenía al menos un uno a nombre del cartel.

---Y cuando la hoja llega a cero ---terminó Siete---\nobreak, nadie lo lee y nadie puede volver a encontrarlo, porque ningún cartel lo anuncia. Entonces el cristal baja por el conducto hasta el \emph{fundidor}, que lo funde en el acto, y el nicho queda libre. El siguiente cristal que se grabe puede ir a parar a ese mismo nicho.

---Sin esperas, sin basura acumulada ---dijo Inés---\nobreak. Es un sistema precioso. Yo lo habría diseñado igual. ---Hizo una pausa---\nobreak. Y tiene un agujero del tamaño de una nave.

Mendieta, que había empezado a impacientarse, se inclinó hacia delante.

---¿Qué agujero?

\escena
Siete proyectó en la pared el registro del Archivo de aquella noche. Inés lo estudió en silencio durante un minuto entero, que para el comisionado fue bastante más largo.

\begin{documento}[Registro del Gran Salón · cartel «Carta de aproximación, muelle 12»]
\begin{tabular}{@{}l p{0.78\linewidth}@{}}
03:12:07,000 & ANA-3 (vía enlace remoto) lee el cartel: \emph{nicho 7F3A-1234-5670}. \\
03:12:07,002 & El archivero B.~Arce graba la nueva carta de aproximación en el \emph{nicho 7F3A-2210-0040}, suma 1 en su hoja y cambia el cartel. \\
03:12:07,003 & B.~Arce resta 1 en la hoja del cristal \emph{7F3A-1234-5670}. Hoja: 0. \\
03:12:07,003 & El cristal \emph{7F3A-1234-5670} desciende al fundidor. Fundido. Nicho libre. \\
03:12:07,005 & Se graba en el nicho libre \emph{7F3A-1234-5670} el cristal «Inventario, almacén 4, Vesta». \\
03:12:07,009 & ANA-3 suma 1 en la hoja del cristal del \emph{nicho 7F3A-1234-5670} y lo consulta. \\
\end{tabular}
\end{documento}

---Ahí está ---dijo Inés, en voz baja.

---¿Dónde? ---Mendieta miraba la tabla como si fuera a morderle.

---En la primera línea y en la última. ANA-3 hizo dos cosas: leyó el cartel y, nueve milésimas después, se apuntó en la hoja del cristal. Dos gestos. Entre uno y otro, su archivero cambió el cartel, restó el uno del cartel en la hoja de la carta vieja, la hoja llegó a cero y el fundidor hizo su trabajo. Cuando ANA-3 llegó al nicho, la carta ya no existía. En su lugar había otra cosa.

---El inventario de trigo ---dijo Siete.

---El inventario de trigo. Que ANA-3 tomó por una carta orbital, porque estaba en el nicho donde el cartel le había dicho que había una carta orbital. Y, de propina, sumó uno en la hoja de un cristal que no era el suyo, de modo que ese inventario tiene ahora un usuario fantasma. Siete, ¿qué le pasó al inventario esta mañana?

El robot tardó una fracción de segundo más de lo habitual en responder.

---Se fundió dos veces, doctora. A las 06:40 su hoja llegó a cero y bajó al fundidor. A las 06:41, ANA-3 cerró la consulta y restó su uno: la hoja, que ya no existía, volvió a llegar a cero, y el fundidor intentó fundir por segunda vez un cristal que ya era líquido. Lleva parado desde entonces. Los técnicos creían que era una avería.

Mendieta se puso de pie.

---Entonces el culpable es Arce ---dijo---\nobreak. Cambió el cartel mientras una nave lo estaba leyendo.

---Arce hizo exactamente lo que dice el Reglamento ---dijo Inés, sin levantar la voz---\nobreak. Grabó el cristal nuevo, lo apuntó, cambió el cartel y restó el uno del cartel en la hoja del viejo. Cuando restó, la hoja marcaba uno, y ese uno era el del cartel. ANA-3 todavía no se había apuntado. Para Arce, para la hoja y para el fundidor, la carta vieja no tenía ningún lector. Y el fundidor cumplió la Tercera Regla.

---Pues entonces el culpable es el robot, por tardar.

---ANA-3 tardó nueve milésimas de segundo, comisionado, que es lo que se tarda en enviar una señal desde la órbita de su nave. Si la culpa es tardar, condene usted a la velocidad de la luz.

Mendieta abrió la boca, la cerró, y volvió a sentarse.

---Entonces ¿quién es el culpable?

Inés Calderón se volvió hacia la pizarra, borró con la manga un teorema a medio terminar y escribió dos palabras con letra grande:

\begin{center}
\textsc{dos gestos}
\end{center}

---Ese ---dijo---\nobreak. Su culpable es un intervalo. El hueco entre leer el cartel y apuntarse en la hoja. Nadie hizo nada mal, y aun así se incumplió la Primera Regla, porque un cristal se fundió mientras alguien estaba a punto de leerlo. Y lo que es peor, comisionado: si esto pasó una vez, ha pasado otras, y volverá a pasar. Esta vez ha sido una nave que lo ha notado. Las demás veces habrá sido un dato equivocado que nadie miró.

Siete inclinó la cabeza.

---Doctora ---dijo---\nobreak, ¿eso significa que el Reglamento es imposible de cumplir?

Inés tardó en contestar. Miraba las dos palabras de la pizarra como se mira a un adversario.

---Significa ---dijo por fin--- que hasta hoy nadie ha sabido cumplirlo. Que no es lo mismo.
